\documentclass{cmspaper}
\usepackage{lineno}
\begin{document}
\begin{linenumbers}

\begin{titlepage}

\cmspas{EXO-09-0XX}
%  \internalnote{2005/000}
%  \conferencereport{2005/000}

%\begin{center}
%{\bf \huge  CMS Physics Analysis Summary }
%\end{center}

\date{\today}

\title{Search for Pair Production of First Generation Scalar Leptoquarks at the CMS Experiment}


  \begin{Authlist}
%   Sarah Eno, Paolo Rumerio, Francesco Santanastasio, Elizabeth Twedt
    The CMS Collaboration, CERN
    %\Instfoot{cern}{University of Maryland, College Park, MD, USA}
  \end{Authlist}


  \begin{abstract}
    The discovery potential for pair production of first generation scalar leptoquarks that 
    decay to an electron and quark is investigated 
    at the CMS experiment using Monte Carlo samples produced with a full simulation of detector response.  
    The analysis strategy assumes an integrated luminosity of 100 pb$^{-1}$ and $pp$ collisions 
    at $\sqrt{s}=10$ TeV.
    Reconstruction and identification of high energy electrons and jets, 
    and optimized selection of events are discussed.
    The electron-jet invariant mass distribution is reconstructed
    and used to identify the possible presence of a leptoquark signal.
    Data-driven techiques are discussed to determine the main standard model backgrounds through use of 
    control samples. The CMS discovery and exclusion potentials using a counting experiment approach 
    are quantified for different leptoquark mass hyphotheses.

%The leptoquarks are assumed to decay exclusively to electrons and quarks.  Detailed explanations of the discovery techniques and data driven methods to extract the signal from background events are given.  
%Both exclusion limits and discovery potential are discussed for masses of the leptoquark ranging from around the current limit of 251 GeV to 1 TeV.
  \end{abstract} 

\end{titlepage}

%=================================================================
\setcounter{page}{2}%JPP

\section{Introduction}
 
This note describes the analysis techniques that will be used to search for evidence of first generation scalar leptoquarks 
at the CMS experiment. The methods described are evaluated using Monte Carlo (MC) simulations assuming a start-up LHC scenario
with about 100 pb$^{-1}$ of integrated luminosity and 5 TeV proton beams.

Leptoquarks (LQ) are new exotic particles conjectured to have both baryon and lepton number and primary decays into leptons and quarks.    
Several models of physics beyond the standard model (SM) predict the existence of leptoquarks, including General Unified Theories, 
technicolor and composite models \cite{theories}.  

The parameters of the model are i) $M_{LQ}$, the LQ mass, ii) $\beta$, the branching fraction for the decay 
$\mbox{LQ} \rightarrow l + q$, where $l$ is a charged lepton and $q$ is a quark, and
iii) $\lambda_{\mbox{LQ}lq}$, the coupling between LQ, lepton and quark ($\mbox{LQ}-l-q$ vertex). 
The complementary decay of a branching ratio 1-$\beta$ describes the decay $\mbox{LQ} \rightarrow \nu + q$, 
where $\nu$ is a weak-isospin partner of the lepton.
Results from HERA on upper limits for LQ production restrict  
$\lambda_{\mbox{LQ}lq}$ to relatively weak coupling\cite{hera}. For example, fixing the parameter 
$\lambda_{\mbox{LQ}lq} = \lambda_{EM} \approx 0.3$, leads to 
a theoretical relative width for the LQ of ($\Gamma_{LQ}/M_{LQ}$) around 0.2\%. 
Consequently, $\Gamma_{LQ}/M_{LQ}$ will not be able to be measured at the LHC, since 
the width of any peak on the mass distribution of the LQ will be dominated by detector resolution.  

Experiments at the Tevatron have placed the most stringent lower limit on the mass of first-generation scalar 
leptoquarks of approximately 290 GeV (assuming $\beta=1$ for the electron channel) \cite{d02008},\cite{cdf2005}.
At the LHC, the pair production of leptoquarks will take place mainly through gluon-gluon fusion and 
quark-antiquark annihilation, yielding $LQ + \overline{LQ}$. 
The cross sections for the two sub-processes are almost independent on the value of 
$\lambda_{\mbox{LQ}lq}$, since there is no $\mbox{LQ}-l-q$ vertex in the Feynman diagram for LQ pair production 
at leading order. At lower rate, the production of single leptoquark in association with a lepton is also possible via quark-gluon 
fusion, yielding $LQ+ l$ or $LQ+ \nu$ (and relative final states with antiparticles). The cross section of single leptoquark production
becomes comparable to the one of pair production only for $M_{LQ}\approx 1$ TeV \cite{LQSingleAndPairProd}, 
which is well above the experimental reach of this start-up analysis (focused on the first 100 pb$^{-1}$ of data at $\sqrt{s}=10$ TeV).  
This note examines the signal for pair production of first-generation scalar leptoquarks that decay into electrons and 
(light) quarks with a branching ratio $\beta=1$. 
For this scenario, the experimental signature of a LQ event is quite striking with two 
high transverse momentum ($p_T$) electrons and two high $p_T$ jets (eejj channel), 
whose combination yields a peak in the electron-jet invariant mass 
spectrum that corresponds to the LQ mass. No true missing transverse momentum is expected in LQ events.


\section{Monte Carlo Samples} \label{sec:MCSamples}
MC signal samples are generated with leptoquark masses ranging from 250 to 600 GeV and $\beta=1$. 
At the LHC, the three main processes that lead to leptoquark pair production are gluon-gluon fusion, 
quark-antiquark annihilation, and the higher order process of quark-gluon fusion. In the mass range to be investigated gluon-gluon fusion dominates. 

The ``Full Simulation'' (FullSim) signal samples are produced using 
the official CMS software for generation (using PYTHIA~\cite{PYTHIA}), 
simulation (based on GEANT4~\cite{GEANT4}), digitization and reconstruction. 
%The above lines are taken from the previous version of the PAS, instead of from the note.
These samples were included in the official ``Summer 08'' production by the CMS Generators group.

``Fast simulation'' (FastSim) signal samples are produced with a slightly later version of the CMS Software.
The FastSim uses a parameterization of the detector response instead of using the full GEANT-based simulation.
The FastSim samples were created during the ``Winter 09'' production campaign by the CMS Generators group.

The SM background is taken from the Summer08 official production samples,
which consists of more than 200M events of various SM processes roughly representing 
the first 100 pb$^{-1}$ of LHC data. The Summer08 samples used for this analysis are listed below:
\begin{itemize}
%
\item $t\bar{t}$+jet events (for N$<$5), generated using MADGRAPH~\cite{MADGRAPH}, inclusive production (all decays); 
%with a LO-NLO K correction factor of 1.85 applied to the LQ cross section;\footnote{The K factor of 1.85 applied 
%to ttbar events is a default in the software provided by CMS for these data samples.  It is from a MCFM NLO vs LO study.}
%
\item $Z/\gamma*$ + N jet events (for N$<$6), generated using MADGRAPH, $Z$ decaying into charged leptons;  
%
\item $W$ + N jet events (for N$<$6), generated using MADGRAPH, $W$ decaying into leptons; 
% $0< P_{T}^{W} < 300 $ GeV (FIXME)
%
%\item QCD multi-jet events, generated with PYTHIA, inclusive production in $P_{T}^{jet}$ bins from 0 to $\infty$;  
\item QCD multi-jet events, generated with MADGRAPH, inclusive production in 
$H_{T}^{jet}$~\footnote{$H_{T}^{jet}=\sum_{jets} p_T^{jet}$} bins from 100~GeV to 1~TeV;  
%
%\item $\gamma$+jet events, generated with PYTHIA, $ 0 < P_{T}^{jet} < 7000 $ GeV;  
%
%\item $WW$, $WZ$, $ZZ$ events, generated with PYTHIA, inclusive production (all decays).
\item $VV$ + jet events (where $VV$ can be $WW$, $WZ$ or $ZZ$), generated with MADGRAPH, $W$ and $Z$ 
decaying into leptons.
\end{itemize} 

Table~1
%\ref{tab:NumEvents} 
lists the signal and background MC samples used in this analysis with their number of events and cross sections.
%shows the number of events and the cross sections at leading order (LO) for samples generated at different LQ mass.

\begin{table}[htb]
  \label{tab:NumEvents}
  \begin{center}
    \begin{tabular}{|l|cccccc|} \hline\hline
%      & \multicolumn{4}{c|}{Leptoquark Mass (GeV)} \\ 
      MC Sample                   & Full/Fast & N. Events & Equivalent             & $\sigma_{NLO}$ (pb) & $\sigma_{LO}$ (pb) & K-factor \\
                                  & Simulation& Analyzed  & Luminosity (pb)        &                     &                    &    \\ 
\hline\hline
      $M_{LQ}=250~$GeV            & Full      & 52k       &    5.15$\times 10^3$   & 10.1                & 6.53               & 1.547\\
      $M_{LQ}=400~$GeV            & Full      & 63k       &      84$\times 10^3$   &  0.75		 & 0.462	      & 1.628\\ 
\hline
      $M_{LQ}=250~$GeV            & Fast      & 125k      &    12.4$\times 10^3$   & 10.1		 & 6.53		      & 1.547\\
      $M_{LQ}=300~$GeV            & Fast      & 127k      &    33.4$\times 10^3$   &  3.8	         & 2.42		      & 1.57\\
      $M_{LQ}=400~$GeV            & Fast      & 150k      &     200$\times 10^3$   &  0.75	         & 0.462	      & 1.628\\
      $M_{LQ}=500~$GeV            & Fast      & 125k      &     635$\times 10^3$   &  0.197  	         & 0.118	      & 1.669\\
      $M_{LQ}=600~$GeV            & Fast      & 131k      &    2.12$\times 10^6$   &  0.0617             & 0.0617             & 1.723\\
%      $M_{LQ}=650~$GeV            & Fast      & 133k      &    3.67$\times 10^6$  &  0.0362 	         & 0.0208	      & 1.740\\
%      $M_{LQ}=700~$GeV            & Fast      & 146k      &    6.70$\times 10^6$  &  0.0218 	         & 0.0124	      & 1.758\\
%      $M_{LQ}=800~$GeV            & Fast      & 146k      &    17.2$\times 10^6$  &  0.0085 	         & 0.00466	      & 1.815\\
%      $M_{LQ}=900~$GeV            & Fast      & 146k      &    41.6$\times 10^6$  &  0.00351	         & 0.00188	      & 1.867\\
%      $M_{LQ}=1000~$GeV           & Fast      & 146k      &    95.4$\times 10^6$  &  0.00153            & 0.000794           & 1.927\\ 
%      $M_{LQ}=1200~$GeV           & Fast      & $>$100k   &                       &  0.000329           & 0.000158           & 2.082\\
%      $M_{LQ}=1500~$GeV           & Fast      & $>$100k   &                       &  0.0000393          & 0.0000165          & 2.382\\ \hline
\hline
      $t\bar{t}$ + N jets         & Full      & 0.905M    &    2.19$\times 10^3$   & 414                 &  317                  & 1.31 \\
      $Z/\gamma$ + N jets         & Full      & 1.159M    &    275	           & 4218                &  3700                 & 1.14 (@14 TeV)\\
      $Z/\gamma$ + N jets         & Fast      & 8.46M     &    2$\times 10^3$      & 4218                &  3700                 & 1.14 (@14 TeV)\\
      $W$ + N jets                & Full      & 6.107M    &     134	           & 45600               &  40000                & 1.14 (@14 TeV)\\
      $VV$ + N jets               & Full      & 101.8k    &    8.63$\times 10^3$   &                     &  11.8                 & \\ \hline
      QCD ($H_T\in[100,250]$~GeV) & Full      & 13.721M   &    0.915	           &                     &  $15.0 \times 10^6$   & \\
      QCD ($H_T\in[250,500]$~GeV) & Full      & 3.698M    &    9.25	           &                     &  $400 \times 10^3$    & \\
      QCD ($H_T\in[500,1000]$~GeV)& Full      & 3.807M    &     272	           &                     &  $14 \times 10^3$     & \\
      QCD ($H_T>1000$~GeV)        & Full      & 0.654M    &    1.77$\times 10^3$   &                     &  $370$                & \\
\hline\hline
    \end{tabular}
    \caption{\small \sl Signal and background MC samples used in this analysis. For each sample, it is indicated 
      if it is generated with the full or fast simulation, the generated number events, 
      and the cross sections at NLO and LO for $\sqrt{s}=10$ TeV. 
      The k-factor is also reported for convenience. Leptoquark cross sections at NLO and LO
      are given based on the calculation discussed at~\cite{Kramer} (thanks to Micheal Kramer 
      for providing the updated numbers for $\sqrt{s}=10$ TeV).}
  \end{center}
\end{table}

\section{Trigger Studies} \label{sec:trig}

The presence of two high energy electrons in the final state is exploited for the online selection 
of candidate leptoquark-antileptoquark events.

Two trigger menus, named ``8E29'' and ``1E31'', have been studied and designed as candidates for 
the two possible luminosity scenario of the CMS start-up.
% More, if needed, at https://twiki.cern.ch/twiki/bin/view/CMS/TSG_18_II_09
In case that the ``8E29'' menu is used, the chosen trigger would be the ``HLT\_Photon15\_L1R''.
The corresponding Level-1 trigger seed is ``L1\_SingleEG8'', which requires an electromagnetic object with 
$E_T>8~$GeV. At the HLT level, the threshold applied is $E_T>15~$GeV and no isolation or track-matching
are required.
If the ``1E31'' menu is used, the above trigger will be prescaled, so the choice would shift to the
``HLT\_Photon25\_L1R'' trigger. The Level-1 trigger seed is still the ``L1\_SingleEG8''. At HLT level the 
threshold is $E_T>25~$GeV with no isolation or track-matching.

The efficiencies for leptoquarks for these triggers have been studied with the FullSim 
samples of Section~\ref{sec:MCSamples} at two leptoquark mass points, and have efficiencies approaching
100\% as shown in Table~\ref{tab:HLTEffic}.

\begin{table}[htbp]
\begin{center}
\begin{tabular}{|c|c|c|c|}
\hline\hline
 $M_{LQ}$ (GeV)     &   ``HLT\_Photon15\_L1R''   &   ``HLT\_Photon25\_L1R'' \\
\hline\hline
250                 & 0.99571  $\pm$ 0.00029     & 0.99343  $\pm$ 0.00035   \\
400                 & 0.99716  $\pm$ 0.00021     & 0.99630  $\pm$ 0.00024   \\
\hline\hline
%% - trigger = HLT_Photon15_L1R
%% eff_LQ250 = 0.99571565802 +/- 0.00028628551
%% eff_LQ400 = 0.99715873016 +/- 0.00021206457
%% - trigger = HLT_Photon25_L1R
%% eff_LQ250 = 0.99342939481 +/- 0.00035412845
%% eff_LQ400 = 0.99630158730 +/- 0.00024184261
\end{tabular}
\end{center}
\caption{Efficiencies of the chosen HLT triggers for LQ masses of 250 and 400~GeV (FullSim samples are used).}
\label{tab:HLTEffic}
\end{table}

\section{Reconstructed Objects} 

\subsection{Electron Studies} \label{sec:electrons}
A reconstructed electron object is formed from deposition of energy in ECAL (called a ``supercluster'') that has an 
extrapolated track pointing to it. This study uses the standard ``GsfElectron'' collection of MC events~\cite{GSFele}.
Standard corrections are applied to set correctly the electromagnetic (e.m.) energy scale.

The criteria used for electron identification (ID) at the reconstruction level are very loose, thereby contributing many false electrons 
to the collection. In this analysis, the HEEP selections for electron ID and isolation~\cite{HEEPNOTE}, which are optimized for 
electrons with energies of hundreds of GeV, have been applied in an effort to: i) keep the efficiency high for electrons 
from LQ decays that have high $p_{T}$, and ii) reduce the number of jets mimicking electrons.
The definitions of the variables used in this analysis are given below:
%
\begin{itemize}
%
\item $H/E$: the ratio of the energy of the closest HCAL Rec Hit to the electromagnetic energy of 
the supercluster associated with the electron; a H/E cut of $<0.1$ is applied at pre-selection level.
%
\item $\sigma_{\eta\eta}$: a variable reflecting the longitudinal distribution of the electron shower (along $\eta$):
\begin{displaymath}
\sigma_{\eta\eta} = \sum( \eta_i - \eta_s )^2 \frac{E_i}{E_{\mbox{s}}} \quad ,
\end{displaymath}
where $E_i$ is the energy in the $i^{th}$ crystal, and $E_s$ is the the energy in the seed cluster of the supercluster associated to the electron.
%
\item $\Delta\eta_{in}$ and $\Delta\phi_{in}$: the difference in $\eta$ and $\phi$ between the track position extrapolated from 
the inner layer tracker to the ECAL surface and the $\eta$ and $\phi$ of the supercluster.
%
\item Number of Tracks ($N_T$): with $p_{T}>1.5$ GeV in an annulus $0.02 < \Delta\mbox{R} < 0.2 $ around the direction of the electron;
%
\item Track Isolation (Track Iso): the sum of the $p_{T}$ of tracks (defined above).
%
%
\item EM Isolation (Iso): the transverse part of the electromagnetic energy 
of all the ECAL Rec Hits (energy deposits in single channels)
%deposited in all the clusters 
in a cone $\Delta\mbox{R} < 0.3$, 
centered on the electron's position in the calorimeter, excluding the Rec Hits that make up the supercluster;
%
\item HAD Isolation: the transverse part of the hadronic energy of all the HCAL Rec Hits in an annulus
$0.15 < \Delta\mbox{R} < 0.3$, centered on the electron's position in the calorimeter. 
%
\end{itemize}

Table \ref{tab:HEEPselection} 
summarizes the cuts from HEEP selections which are 
used in this analysis. 

\begin{table}[htbp]
  \label{tab:HEEPselection}
  \begin{center}
    \begin{tabular}{|lcc|lcc|} \hline
      \multicolumn{3}{|c|}{ID Variables} & \multicolumn{3}{|c|}{Isolation Variables} \\ 
      Variable & Barrel & Endcap & Variable & Barrel & Endcap  \\ \hline
      $H/E$  & $<0.05$ & $<0.1$ & $N_T$  & $<4$ & $<4$ \\ \hline
      $\sigma_{\eta\eta}$  & $<0.011$ & $<0.0275$ & Track Iso (GeV) & $<7.5$ & $<15$ \\ \hline
      $|\Delta\eta_{in}|$  & $<0.005$ & $<0.007$ & EM Iso (GeV) & $<6+0.01*E_{t}$ & $<6+0.01*E_{t}$ \\ \hline
      $|\Delta\phi_{in}|$  & $<0.09$ & $<0.09$ & HAD Iso (GeV) & $<4+0.005*E_{t}$ & $<4+0.005*E_{t}$ \\ \hline
    \end{tabular}
  \caption{\small \sl HEEP electron ID and isolation criteria. Reconstructed electrons are associated to the 
    barrel (endcaps) if $|\eta|<1.442$ ($1.560<|\eta|<2.5$).}
  \end{center}
\end{table}

Electrons from LQ decays, being usually isolated from jets and having a $p_{T}$ of 
hundreds of GeV, can be reconstructed with high efficiency. 
Table \ref{tab:ElecEffAcc} 
shows the electron acceptance $A_{gen}$ (the acceptance cuts are $P_{T}^{gen}>$ 30 GeV 
and $|\eta^{gen}|<2.5$), the reconstruction efficiency $\varepsilon_{reco}$ within the acceptance 
region, and the overall reconstruction efficiency $A_{gen} \times \varepsilon_{reco}$ 
for single electrons from LQ decays for a 400 GeV LQ mass.
The overall reconstruction efficiency on a single electron, including acceptance, is $\approx 80\%$.
\begin{table}[htb]
  \label{tab:ElecEffAcc}
  \begin{center}
    \begin{tabular}{|l|c|c|c|} \hline
      LQ mass (GeV) & $A_{gen}$ & $\varepsilon_{reco}$ & $A_{gen} \times \varepsilon_{reco}$\\ \hline
      400 & 96.4\% & 83.4\% & 80.4\% \\ \hline
    \end{tabular}
    \caption{\small \sl Electron acceptance ($A_{gen}$), 
      reconstruction efficiency after ID/Isolation cuts ($\varepsilon_{reco}$), and overall reconstruction efficiency 
      ($A_{gen} \times \varepsilon_{reco}$) for single electrons from decays of LQ with mass of 400 GeV (FullSim sample used).   
      Relative statistical uncertainties on efficiencies are less than 0.2\%.  
      } 
  \end{center}
\end{table}

\subsection{Jet Studies} \label{sec:jet}

The jet reconstruction algorithm used in this analysis is an interactive cone algorithm using a radius of $\Delta R=0.5$ \cite{JetAlg}.  
In order to remove the electron contamination from the jet collection (disambiguation)
the collection of electrons that pass electron ID and Isolation 
selection criteria is created, as described in Section~\ref{sec:electrons}. The three-vectors of the 
electrons in this new collection is then compared to the three-vectors of the jets. 
A jet is removed if it is within a $\Delta R$ of 0.5 from one of the electrons. 


The corrections for the Jet Energy Scale (JES) are divided into various levels in the CMS Software \cite{JES}.  
In this analysis the jet energy is corrected with the level 2, 3, and 5 corrections.

The Level 2 jet energy corrections comprise the 
``relative jet corrections'' that provide uniform jet energy response in $\eta$.  
The Level 3 corrections are the ``absolute jet corrections'', that aim to correct  the jet energies in the 
region of $\eta < 1.3$ back to the particle level energy, to assure agreement in energy with the generated MC samples. 
The relative and absolute corrections are the standard CMS Software corrections recommended for all jet analyzes,
and on average change the energies by 20\%.
In addition, Level 5 corrections are applied to take account of the flavor of the originating parton.
The flavor correction on average changes the energy of the jets by only a few percent.

\section{Event Selection} \label{sec:eventSelection}
The basic strategy to identify the existence of a particle that decays to a jet and an electron 
is to study the invariant mass of the electron-jet pairs in the events. 
The signal of a leptoquark would appear as a bump in the distribution of this invariant mass.
The following selections have been made in order to minimize the number of background events 
selected while retaining a high signal efficiency.

The online selection of candidate events is made by the HLT triggers EMHighEt and EMVeryHighEt 
described in Section~\ref{sec:trig} while the offline identification of electrons and jets 
proceeds as described in Sections~\ref{sec:electrons} and \ref{sec:jet}, respectively.
The offline selection of the eejj sample continues with the following cuts:

\begin{enumerate}
\item at least 2 isolated electrons, one required to have $P_T>85$~GeV, 
and the other $P_T>30$~GeV 
\item at least 2 jets, both required to have $P_T>50$~GeV
\item $M_{ee}>100$~GeV
\item $S_T\equiv P_T(e_1)+P_T(e_2)+P_T(j_1)+P_T(j_2)>400$~GeV
\end{enumerate}

The $P_T$ threshold for one electron in cut 1 is set to 5~GeV above the threshold of the EMHighEt trigger 
in order to increase the confidence that the trigger was actually fired by a real electron in all events.
The other electron will have a $P_T>30$~GeV as it is implicit in the definition of electron. 
Cut 3, where $M_{ee}$ is the invariant mass of the electron pair, removes background events from 
$Z/\gamma$+jets events as shown in Figure~\ref{fig:Mee_St_distributions}-left.
Cut 4, where the variable $S_T$ is defined as the scalar sum of the transverse momenta of the 
2 electrons and 2 jets with highest $P_{T}$ of the event, 
is applied following the approach of the experiment $D0$ in 
\cite{Abazov:2001mx}. In that paper, an event selection optimization as a function of
the combinations of several kinematic variables has shown that $S_T$ is the most powerful one 
and little is gained by adding cuts on other variables. The distribution of $S_T$ for the present
analysis is shown in Figure~\ref{fig:Mee_St_distributions}-right.


Once the eejj sample is built, there are two ways to combine two electrons and two jets to make two electron-jet pairs. 
For each event, the combination with the minimum difference, $\Delta M_{ej}$, between the invariant masses, $M_{ej}$, 
of the two electron-jet pairs is chosen. 
The resulting $M_{ej}$ distribution is shown in Figure~\ref{fig:Mej_allComb} for the signal and remaining backgrounds. 
A detailed study of the optimization of the selection criteria, and of alternative algorithms to identify the best candidate 
electron-jet combinations will be performed in future upgrades of this analysis.

\begin{figure}[htbp]
  \begin{center}
    \begin{tabular}{cc}
      \resizebox{7.5cm}{!}{\includegraphics{plots/UMD.eps}} &
      \resizebox{7.5cm}{!}{\includegraphics{plots/UMD.eps}} \\
    \end{tabular}
    \caption{\small \sl Left: invariant mass of the electron pair, $M_{ee}$. 
             Right: scalar sum of the $P_T$ of the 2 leading electrons and 2 leading jets. 
	     In each histogram, the distributions for the signal (for leptoquark mass 250 and 650~GeV) and the 
	     relevant backgrounds are shown after applying all cuts except the one involving the 
	     plotted variable. 
	     The background histograms are summed on top of each other.
	     The histogram labeled with ``Others'' contains: 1) for $M_{ee}$, $W$+jets (33\%), 
	     $WW$ (11\%), $ZZ$ (20\%), and $WZ$ (36\%) events; 2) for $S_{T}$, $W$+jets (57\%), 
	     $WW$ (19\%), $ZZ$ (8\%), $WZ$ (14\%), and $\gamma$+jet (2\%) events.
	     With the available MC statistics, no QCD multi-jet events pass 
	     the selection.}
	     %Arrows indicate the cut subsequently applied.
    \label{fig:Mee_St_distributions}
  \end{center}
\end{figure}


\begin{figure}[htbp]
  \begin{center}
    \resizebox{10cm}{!}{\includegraphics{plots/UMD.eps}}
    \caption{\small \sl Distribution of the invariant mass, $M_{ej}$, of the electron-jet pairs 
      with smaller $\Delta M_{ej}$
      for signal (at 250~GeV and 650~GeV LQ mass) and backgrounds. 
      The complete event selection has been applied.
      The background histograms are summed on top of each other.
      The histogram labeled with ``Others'' contains $W$+jets (65\%), 
      $WW$ (20\%), $ZZ$ (8\%), and $WZ$ (7\%) events. 
      With the available MC statistics, no QCD multi-jet and $\gamma$+jet events pass 
      the selection.}
    \label{fig:Mej_allComb}
  \end{center}
\end{figure}

The efficiency of each selection cut is shown in Table \ref{tab:selection_effic_250} 
for a LQ sample with a mass of 250 GeV.  
Table~\ref{tab:selection_effic_ttbar} shows the number of events selected by each cut 
for $t\bar{t}$ and $Z/\gamma$+jet events, which are the dominant backgrounds in the eejj sample. 
A summary of the number of selected signal and background events expected in 100 pb$^{-1}$ of data 
is reported in Table \ref{tab:EventSelSummary}. 
The overall signal selection efficiencies 
are around 30-50\% for the LQ masses investigated. 
After the described event selection, the dominant background 
contributions are $t\bar{t}$ and $Z/\gamma$+jets. 

\begin{table}[htbp]
\begin{center}
\begin{tabular}{|c|c|c|}
\hline
\hline
 & $N_{ev}$ $100pb^{-1}$ & $\varepsilon$ \\
\hline
\hline

no cut &1400 $\pm$ 44& - \\
+ hlt (HE or VHE) &1121 $\pm$ 40 & 8.01e-01 $\pm$ 1.26e-02\\
+ 2 ele (ID) $P_{T} >30$ GeV &644 $\pm$ 30 & 5.74e-01 $\pm$ 3.36e-02\\
+ 2 ele (ID+Iso) $P_{T} >30$ GeV &592 $\pm$ 29 & 9.20e-01 $\pm$ 6.19e-02\\
+ 2e, 2j $P_{T}^{ele}$ ($P_{T}^{jet}$) $>$30(50) GeV &508 $\pm$ 27& 8.58e-01 $\pm$ 6.14e-02\\
+ $P_{T}$ $1^{st}$ ele $>$ 85 GeV &506 $\pm$ 27& 9.97e-01 $\pm$ 7.41e-02\\
+ $M_{ee} >100$ GeV& 455 $\pm$ 25& 8.98e-01 $\pm$ 6.86e-02\\
+ $S_{T} >400$ GeV &425 $\pm$ 24& 9.35e-01 $\pm$ 7.46e-02\\
\hline

full selection efficiency& & 3.04e-01 $\pm$ 1.45e-02\\
\hline
\end{tabular}
\end{center}
\caption{\small \sl LQ sample with $M_{LQ}=250$ GeV: the first column lists the selection sequence, $N_{ev}$ $100pb^{-1}$ is the number of selected events in $100pb^{-1}$, $\varepsilon$ is the relative selection efficiency with respect to the previous cut. The full selection efficiency is shown in the last line.}
\label{tab:selection_effic_250}
\end{table}



\begin{table}[htbp]
\begin{center}
\begin{tabular}{|c| |c|c|}
\hline
\hline
 & $t\bar{t}$ sample  & $Z/\gamma$ sample\\
 & $N_{ev}$ $100pb^{-1}$ & $N_{ev}$ $100pb^{-1}$ \\
  
\hline
\hline

hlt (HE or VHE) &2115.9 $\pm$ 9.6 & 2466 $\pm$ 14 \\
+ 2 ele (no Iso) pT $>30.0$ GeV &137.4 $\pm$ 2.5& 1356 $\pm$ 12 \\
+ 2 ele (Iso) pT $>30.0$ GeV &105.4 $\pm$ 2.1& 1329 $\pm$ 12  \\
+ 2e, 2j pT$_{ele}$(pT$_{jet}$) $>$30.0(50.0) GeV &49.6 $\pm$ 1.5& 124.7 $\pm$ 3.0\\
+ $P_{T}$ $1^{st}$ ele $>$ 85 GeV &47.2 $\pm$ 1.4& 115.7 $\pm$ 2.9\\
+ $Mee >100.0$ GeV&40.2 $\pm$ 1.3& 10.82 $\pm$ 0.89 \\
+ $St >400.0$ GeV &21.26 $\pm$ 0.98 & 5.10 $\pm$ 0.60 \\
\hline
\end{tabular}
\end{center}
\caption{\small \sl $t\bar{t}$ and $Z/\gamma$ samples: the first column lists the selection sequence, $N_{ev}$ $100pb^{-1}$ is the number of selected events in $100pb^{-1}$.}
\label{tab:selection_effic_ttbar}
\end{table}



%\begin{table}[htbp]
%\begin{center}
%\begin{tabular}{|c| |cc| |cc|}
%\hline
%\hline
% & $t\bar{t}$ sample &              & $Z/\gamma$ sample &\\
% & $N_{ev}$ $100pb^{-1}$ & $N_{ev}$ &  $N_{ev}$ $100pb^{-1}$ & $N_{ev}$ \\
%  
%\hline
%\hline
%
%hlt (HE or VHE) &2115.9 $\pm$ 9.6& 48359.0 & 2466 $\pm$ 14 & 30413 \\
%+ 2 ele (no Iso) pT $>30.0$ GeV &137.4 $\pm$ 2.5& 3139.0 & 1356 $\pm$ 12 & 13136 \\
%+ 2 ele (Iso) pT $>30.0$ GeV &105.4 $\pm$ 2.1& 2408.0 & 1329 $\pm$ 12 & 12830 \\
%+ 2e, 2j pT$_{ele}$(pT$_{jet}$) $>$30.0(50.0) GeV &49.6 $\pm$ 1.5& 1117.0 & 124.7 $\pm$ 3.0& 1807 \\
%+ $P_{T}$ $1^{st}$ ele $>$ 85 GeV &47.2 $\pm$ 1.4& 1063.0 & 115.7 $\pm$ 2.9& 1694 \\
%+ $Mee >100.0$ GeV&40.2 $\pm$ 1.3& 1013.0 & 10.82 $\pm$ 0.89 & 157 \\
%+ $St >400.0$ GeV &21.26 $\pm$ 0.98 & 476 & 5.10 $\pm$ 0.60 & 76 \\
%\hline
%\end{tabular}
%\end{center}
%\caption{\small \sl $t\bar{t}$ and $Z/\gamma$ samples: the first column lists the selection sequence, $N_{ev}$ $100pb^{-1}$ is the number of selected events in $100pb^{-1}$, $N_{ev}$ is the absolute number 
%of selected events (unweighted for the cross section).}
%\label{tab:selection_effic_ttbar}
%\end{table}
%

\begin{table}[htbp]
\begin{center}
\begin{tabular}{|cccc||ccc|}
\hline
                 &   Signal&         &          & & Background & \\
$M_{LQ}$=250 GeV & 400 GeV & 650 GeV & 1000 GeV & $t\bar{t}$  & $Z/\gamma$+jet & Others \\
\hline
425 $\pm$ 24 & 52.4 $\pm$ 2.5 & 3.4 $\pm$ 0.1 & 0.21 $\pm$ 0.01 & 21.3 $\pm$ 1.0 & 5.1 $\pm$ 0.6 & 1.1 $\pm$ 0.3 \\
\hline
\end{tabular}
\end{center}
\caption{\small \sl Number of selected events in the eejj sample expected in 100 pb$^{-1}$ of data for signal events at different LQ masses 
and for background events. 
The LQ cross section rapidly falls at high LQ mass, thus 
producing a relative decrease in the number of selected events. 
In the background part, ``Others'' includes $W$+jet, and 
di-boson ($WW$,$WZ$,$ZZ$) events. 
With the available MC statistics, no QCD multi-jet and $\gamma$+jet events pass the final selection.}
\label{tab:EventSelSummary}
\end{table}


\section{Data-driven techniques for background estimate} \label{sec:bkgStudy}
This section describes data-driven techniques used to estimate the shape of the $M_{ej}$ distribution for $t\bar{t}$ and $Z/\gamma$+jet backgrounds using 
control samples. These shapes are used in a fit to the data to extract the number of signal events, as described 
in Section~\ref{sec:signalExtraction}. The methods for background estimation described here will not rely on the Monte Carlo 
once real data is available. They are
therefore particularly suited for first data taking when the confidence that the MC describes well the data is expected to be limited.
\subsection{$t\bar{t}$ background shape}
A good control sample can be obtained by using the same selection criteria applied for the eejj sample, but 
requiring at least one electron and one muon (e$\mu$jj sample) instead of 2 electrons in the final state, in addition to the two jets. 
For $t\bar{t}$ events, the $M_{lj}$ distributions of the e$\mu$jj and the eejj samples
are expected to be very similar in shape since the kinematics of the process does not depend 
on the nature of the lepton. 
 Figure~\ref{fig:ttbar} shows a good agreement between 
the shape of the two distributions with the current MC statistics available. 
With 100 pb$^{-1}$ of data the e$\mu$jj sample is expected to have about 30 events ($\approx 60$ entries in the 
$M_{lj}$ distribution). 
The e$\mu$jj sample is enriched in $t\bar{t}$ events with a small contamination of about 
2\% of $WW$ and $Z/\gamma$+jet events. 

%For a sample of $t\bar{t}$ events at generator level, the number of e$\mu$jj events is expected to be exactly two times the number of 
%eejj events, considering all the possible combinations of the $W$ decays. The trigger filter, the offline selection, 
%and the different reconstruction efficiency, acceptance and $P_{T}$ resolution between electron and muon
%put a bias in the relative amount of eejj and e$\mu$jj events.
%In general it is possible to correct for these effects in order to estimate the number of eejj in the signal sample directly from 
%the size of the e$\mu$jj control sample. In this way the control sample can be used to estimate both the normalization and the shape
%of the $t\bar{t}$ background. In this analysis only the shape of the e$\mu$jj sample is used in the signal extraction, but
%the possibility to get also the normalization from the control sample will be investigated in future upgrades of this analysis.

\begin{figure}[htb]
  \begin{center}
  \begin{tabular}{cc}
  \resizebox{10cm}{!}{\includegraphics{plots/UMD.eps}} \\ 
  \end{tabular}
  \caption{Distributions of the lepton-jet invariant mass for the eejj and the e$\mu$ samples, for $t\bar{t}$ events}
  \label{fig:ttbar}
  \end{center}
\end{figure}

\subsection{$Z/\gamma$+jet background shape} 

The same concept of control sample can be applied to estimate the $M_{ej}$ shape of the $Z/\gamma$+jet background.
In this case the control sample can be obtained by using the same selection criteria applied for the eejj sample except the $M_{ee}$ cut, which is modified to select events with a real $Z$ boson reconstructed ($ 80\mbox{ GeV} < M_{ee} < 100\mbox{ GeV}$). 
This control sample is an almost pure sample of  
$Z/\gamma$+jet events (about 4\% contamination dominated by $t\bar{t}$ and $WZ$ events) since the cross section of the process is resonant at the Z mass, and is 
independent from the signal sample by construction. In addition the control sample has about 10 times the number of $Z/\gamma$+jet events in 
the signal sample (about 50 events expected in the control sample for 100 pb$-1$ of data), which 
is important to reduce the statistical errors in the determination of the background shape.

Figure~\ref{fig:zjet} shows the good agreement between the shape of the $M_{ej}$ distribution for the two samples 
with the current MC statistics available. The agreement observed can be interpreted as the consequence of the 
weak correlation between the two reconstructed quantities $M_{ej}$ and $M_{ee}$ for selected events. Indeed, it is observed that the selected electrons acquire a high $P_{T}$ due to
the substantial boost, rather than a high mass, of the $Z/\gamma$ bosons.  

\begin{figure}[htb]
  \begin{center}
  \begin{tabular}{cc}
  \resizebox{10cm}{!}{\includegraphics{plots/UMD.eps}} \\ 
  \end{tabular}
  \caption{Distributions of the electron-jet invariant mass for the signal eejj sample 
    (cut 1,2,3,4)
    and the control sample (cut 1,2,4 + $80\mbox{ GeV} < M_{ee} < 100\mbox{ GeV}$), for $Z/\gamma$+jet events.}
  \label{fig:zjet}
  \end{center}
\end{figure}

\section{Signal Extraction} \label{sec:signalExtraction}

This section describes the procedure and verifies the ability to extract the 
number of signal events from a distribution of the invariant mass, $M_{ej}$, of the 
electron-jet pairs with the characteristics expected from the studies 
of the previous sections.

MC experiments are generated to produce $M_{ej}$ distributions that are the sum
of a leptoquark signal and the only relevant backgrounds: $t\bar{t}$ and $Z/\gamma$+jets.
The average leptoquark signal distribution, $h_s$, is taken from the FullSim MC 
sample at the leptoquark mass under consideration after applying the selection
described in Section~\ref{sec:eventSelection}.
The average distributions, $h_{t\bar{t}}$ and $h_{Z/\gamma\mathrm{+jets}}$ 
of the $t\bar{t}$ and $Z/\gamma$+jets backgrounds are determined from the control samples
and procedures described in Section~\ref{sec:bkgStudy}.

A log-likelihood fit is performed on the histogram of each generated MC experiment
using a function
\begin{displaymath}
  h_{s+b} = N_s \cdot \tilde{h}_s + N_b \cdot \tilde{h}_b
\end{displaymath}
where $N_s$ and $N_b$ are the free parameters of the fit for the number of signal and background events,
\begin{displaymath}
  \tilde{h}_b = \frac{R}{1+R} \tilde{h}_{t\bar{t}} + \frac{1}{1+R} \cdot \tilde{h}_{Z/\gamma\mathrm{+jets}}~\mathrm{,}
\end{displaymath}
and $\tilde{h}_x$ represents a general distribution $h_x$ normalized to unity. 

The ratio
\begin{displaymath}
  R \equiv N_{t\bar{t}} / N_{Z/\gamma\mathrm{+jets}}~\mathrm{,}
\end{displaymath}
between the number of $t\bar{t}$ and $Z/\gamma$+jets events is currently determined from 
MC and has a value $R_{MC}=4.2$.
Figure~\ref{fig:Mej_fit} shows a MC experiment and the result of the fit.
%A tendency of the fit to underestimate the number of signal events is observed. 
%Such an effect amounts to about 10\% with a moderate dependency on the leptoquark mass and the 
%integrated luminosity. This effect may bias the results towards an overestimation of the integrate%d luminosity needed for discovery.

The fitting procedure described in this section is used in the following section
to determine the CMS potential for discovering a leptoquark signal. 


 \begin{figure}[htb]
   \begin{center}
     \resizebox{10cm}{!}{\includegraphics{plots/UMD.eps}}
     \caption{\small \sl A MC experiment (dots with error bars) generated for a leptoquark mass of 
       400~GeV (blue-filled histogram), a sum of the $t\bar{t}$ and $Z/\gamma$+jets 
       backgrounds (gray-filled histogram) with a relative ratio R=4.2 determined from 
       full simulation MC, and a number of signal events $N_s=50$ corresponding
       to an integrated luminosity of 95~$pb^{-1}$.
       The blue-line open histogram is the result of the fit and the estimate of signal 
       events is 48.5$\pm$ 3.8.}
     \label{fig:Mej_fit}
   \end{center}
 \end{figure}



\section{CMS Potential} \label{CMSpotential}

In order to estimate the integrated luminosity needed by CMS to claim 
a discovery in the electron channel of the first generation leptoquark, 
several MC experiments with signal and background events are produced
as described in Section~\ref{sec:signalExtraction}.

The significance of each experiment is extracted using the log-likelihood estimator 
\begin{displaymath}
S_L \equiv \sqrt{2\cdot ln{(L_{s+b}/L_{b})}}~\mathrm{,}
\end{displaymath}
where $L_{s+b}$ and $L_b$ are the maximum likelihood values for the fit to
the signal plus background and background only hypothesis, respectively.

A set of MC experiments is generated for a leptoquark mass and a value of 
the integrated luminosity.
The distribution of $S_L$ of the set of experiments is approximately gaussian 
and its median is used as the estimate of the significance.
The procedure is repeated for leptoquark masses of 250, 400 and 650~GeV and
different values of the integrated luminosity. 
 The results are shown in
Figure~\ref{fig:sign_vs_Lint_sysR}, which includes also the effect of the systematic 
uncertainties discussed in the next section. 

\subsection{Systematic uncertainties}

%The impact of systematic effects on the discovery potential and the determination 
%of the experiment sensitivity are not yet included in this study.

%A study for determining the sensitivity of the experiment will be performed.
%It will be used to set upper limits in case a signal is not observed. 
%It will be used an approach similar to the one used in \ref{highmassToMuons}.
%The same tools to generate MC experiments and perform log-likelihood fits will be used.

%Systematics: Luminosity(+_10\%), JES (+_10\%), ratio ttbar/Zgammajets, etc.

Once real data is available, and given the procedure used to extract the signal described in 
Section~\ref{sec:signalExtraction},
only the uncertainties that affect the knowledge of the shape of the distributions of 
the total background and the signal will have an impact on the significance, and have to be 
accounted for as systematic uncertainties.

The value of $R$, defined in Section~\ref{sec:signalExtraction}, 
directly affects the shape of the total background distribution used in the fit by changing the
mixture of the two backgrounds. 
The current knowledge, $R_{MC}$, of $R$ comes from the Monte Carlo
study and may introduce a deviation $\delta \equiv R/R_{MC}$ from the true value of $R$
due, for example, to imperfect knowledge of the relative magnitude of the two background 
cross sections and respective selection efficiencies.

$\tilde{R}$ is defined as the ratio between the number of $t\bar{t}$ and $Z/\gamma$+jets events in the
control samples described in Section~\ref{sec:bkgStudy}, and $\tilde{R}_{MC}$ is its
estimate from Monte Carlo. It is expected that the main reasons that
make $\delta$ deviate from unity affect $\tilde{R}/\tilde{R}_{MC}$ in the same way, assuming that the MC uncertainties are
dominated by the limited knowledge of the cross sections.
Therefore, once real data is available to calculate $\tilde{R}$ from the control samples of $t\bar{t}$ and $Z/\gamma$+jets, 
the estimate $\delta \simeq \tilde{R}/\tilde{R}_{MC}$ 
will be used to improve the knowledge of $R$ from the current value, $R_{MC}$, to $R_{data+MC} \equiv \delta \cdot R_{MC}$.

The uncertainty on the estimate of $\delta$ given by $\tilde{R}/\tilde{R}_{MC}$ will be 
dominated by the statistical error of $\tilde{R}$ from the number of events in the control samples, and will propagate to $R_{data+MC}$. 
To estimate the systematic effect on the significance, the fit is repeated by varying the 
value of $R$ by one standard deviation from its central value 
$R_{MC}$ \footnote{Not having real data yet, 
exact correctness of the MC is currently assumed by setting $\delta=1$, 
making effectively $R_{data+MC}=R_{MC}$.}, as shown in 
Figure~\ref{fig:sign_vs_Lint_sysR}.
%The systematic uncertainty due to $R$ is larger, for the same integrated luminosity, for a
%LQ mass of 250~GeV. This is due to the topology of the $M_{ej}$ distribution 
%of the signal that, at small LQ masses, peaks where the background does.
%For higher values of the LQ mass, the significance determination seems to be little 
%sensible to variations of $R$. This may be understood by considering the increased separation
%between the main features of the signal and backgound distributions, and the relative similarity 
%between the two backgound shapes: 
%the mixture between $t\bar{t}$ and $Z/\gamma$+jets events may vary, but
%the mixture between total background and signal does little.
As an example, for the point at $M_{LQ}=400$~GeV, integrated luminosity=100~$pb^{-1}$ and significance=12.7, 
the uncertainty on $R_{data+MC}$ is about 15\%, but this determines a change of 
significance of only 1\%.

The uncertainty 
of $R$ tends to be ``absorbed'' in the fitting procedure and have a limited impact on the 
significance determination.
In addition, this reassures that once the real data is available and the value of the correction 
factor $\delta$ can be estimated as described above,
even a significant departure from unity will not require a drastic re-assessment of
the current estimate of the luminosity needed for discovery.
The uncertainty on the knowledge of the signal shape used in the fit may also produce a systematic
uncertainty on the significance. A study of this effect will have to be performed.

In the case of setting an upper limit in the case of absence of signal, a larger number of systematic
uncertainties will have to be considered (on the integrated luminosity, the detector efficiency and resolution, the jet energy scale, and theoretical uncertainties) 
as they do have a direct impact on it. 
The procedure for setting an upper limit and the study of the 
systematics involved is not yet included in this version of the analysis.


 \begin{figure}
   \begin{center}
     \resizebox{13cm}{!}{\includegraphics{plots/UMD.eps}}
     \caption{Significance of the discovery as a function of the integrated luminosity.
       For each LQ mass, the fit to extract the significance is repeated varying the value of $R$
       (the ratio between the number of $t\bar{t}$ and $Z/\gamma$+jets events in the final sample)
       of one standard deviation from its central value. 
       The horizontal line represents a 5 sigma discovery.}
     \label{fig:sign_vs_Lint_sysR}
   \end{center}
 \end{figure}


\section{Conclusion}

A MC-based analysis of the pair production of first generation leptoquarks that decay to
an electron and a jet has been presented.
The analysis strategy and detector performance description corresponds to a 100$~pb^{-1}$ scenario.
Standard CMS techniques are used for electron and jet identification. 
A cut-based event selection has been applied and its optimization study will come in a future upgrade
of the analysis.
Several background sources have been studied and only two of them provide a significant contribution
after event selection. 
Data-driven techniques to understand the characteristics of these contributions have been developed 
and studied.

The signal component is extracted with a log-likelihood fit that does not rely on precise knowledge
of the level of the background but rather by exploiting the different shapes of the signal and background
distributions. 
The discovery potential of this analysis has been determined using a likelihood 
ratio estimator. This has shown that, for leptoquark masses just above the Tevatron exclusion limit of 290~GeV
\cite{d02008}, an early discovery is possible with a few $pb^{-1}$.
With an integrated luminosity of 100~$pb^{-1}$, discovery should be possible up
to LQ mass of about 500 GeV.
The impact of the systematics on the discovery potential has been discussed.
The procedure for setting an upper limit in absence of signal will be studied 
in the next upgrade of the analysis.


%--------------------------------------
%\clearpage
%--------------------------------------

\begin{thebibliography}{}

\bibitem {theories} {D.Acosta and S.K.Blessing, Ann.Rev.Nucl.Part.Sci 49,389},
  1999,
  {\em From June 2005}
\bibitem{hera}{H1 collaboration, Search for leptoquark bosons in ep collisions at HERA}, hep-ex/0506044,
\bibitem{d02008}{Search for First-Generation Leptoquarks in the dielectron channel with the DO Detector in $p\bar{p}$ Collisions at $\sqrt{s}=1.96$ TeV}, Jun 2008,
  {\em DO Note 5644-CONF}
\bibitem{cdf2005}{Search for first-generation scalar leptoquarks in $p\bar{p}$ collisions at $\sqrt{s}=1.96$ TeV}, Sept 2005,
  {\em Phys. Rev. D 72, 051107 (2005)} (CDF paper)
\bibitem{LQSingleAndPairProd}{Leptoquark Single and Pair production at LHC with CalcHEP/CompHEP in the complete model}, Feb 2005,
  {\em arXiv:hep-ph/0502067, JHEP0509:005,2005, MSU-HEP-070204}
  
  %  \bibitem{d02007}{Searches for LQ production at D0}, Oct 2007,
  %    {\em arXiv0710.0255v1}
  
  %\cite{Mangano:2002ea}
\bibitem{MADGRAPH}
  J. Alwall et Al.
  ``MadGraph/MadEvent v4: The New Web Generation''
  [arXiv:0706.2334].
  
  %  \bibitem{Mangano:2002ea}
  %    M.~L.~Mangano, M.~Moretti, F.~Piccinini, R.~Pittau and A.~D.~Polosa,
  %    ``ALPGEN, a generator for hard multiparton processes in hadronic collisions,''
  %    JHEP {\bf 0307} (2003) 001
  %    [arXiv:hep-ph/0206293].
  %    %%CITATION = JHEPA,0307,001;%%
  
\bibitem{HeepHlt}{D. Acosta et al., The CMS High Level Trigger}, CMS AN 2007/009,
\bibitem{Kramer}{M.Kramer et al., Pair production of scalar leptoquarks at the LHC},Jan 2008 ,{\em arXiv 0411038v2}
  %\cite{Abazov:2001mx} 	 
\bibitem{GSFele}{S. Baffioni et al., Electron Reconstruction in CMS}, CMS NOTE 2006/040,
\bibitem{HEEPNOTE} {CMS Collaboration, ``Search for high mass resonance production decaying into an electron pair in the CMS experiment''}, CMS PAS EXO-08-001
\bibitem{EleID}{D. Newbold et al., Electron ID at High Energies}, CMS AN 2008/045,
\bibitem{TagAndProbe}{CMS collaboration, Measuring Electron Efficiencies at CMS with Early Data}, CMS PAS EGM-07-001 (2007),
\bibitem{JetAlg} {P. Schieferdecker et al., Performance of Jet Algorithms in CMS}, CMS AN-2008/01,
\bibitem{Abazov:2001mx} 	 
  V.~M.~Abazov {\it et al.}  [D0 Collaboration], 	 
  %``Search for first-generation scalar and vector leptoquarks,'' 	 
  Phys.\ Rev.\  D {\bf 64} (2001) 092004 	 
  [arXiv:hep-ex/0105072]. 	 
  %%CITATION = PHRVA,D64,092004;%%
  %  \bibitem{highmassToMuons}{\bf I. Altsybeev et al., Search for new high-mass resonances decaying to muon pairs in the CMS experiment}, CMS AN 2007/038,
\bibitem{JES} {R. Harris, K. Kousouris, MC Truth L2 \& L3 Factorized Jet Corrections at CMS}, CMS AN-2008/03,
\bibitem{PTDR} {CMS Collaboration 2007, CMS Technical Design Report, Volum II: Physics Performance} 
J. Phys. G.: Nucl. Part. Phys. 34 995-1579, ``Luminosity Uncertainty'' at page 1500
\bibitem{PDFRescaling} {Parton Distribution Uncertainty Determination within CMSSW - CMS AN-2009/048}, FIXME
\bibitem{ref:scp} S. Bityukov, N. Krasnikov, S. Erofeeva, and A. Nikitenko, Program for evaluation of significance,
  confidence intervals and limits by direct calculation of probabilities, PhyStat 2005, Oxford, UK 
  [http://cmsdoc.cern.ch/~bityukov]
\bibitem{ref:bayes}  C. Amsler et al., Physics Letters B667, 1 (2008), section 32.3.1
  
\end{thebibliography}

\end{linenumbers}
\end{document}
